\documentclass{article}
\usepackage{amsmath, amssymb, amsthm}
\usepackage{hyperref}
\usepackage{geometry}
\geometry{margin=1in}

\title{Erd\H{o}s Problem \#948}
\date{Last updated: 2025-08-31}
\author{Source: erdosproblems.com}

\begin{document}
\maketitle

\noindent\textbf{Status:} OPEN \quad \textbf{No prize}

\noindent\textbf{Tags:} number theory, ramsey theory

\medskip

\noindent\textbf{Problem Statement:}

Is there a function $f(n)$ and a $k$ such that in any $k$-colouring of the integers there exists a sequence $a_1&#60;\cdots$ such that $a_n&#60;f(n)$ for infinitely many $n$ and the set\[\left\{ \sum_{i\in S}a_i : \textrm{finite }S\right\}\]does not contain all colours?




Erd\H{o}s initially asked whether this is possible with the set being monochromatic, but Galvin showed that this is not always possible, considering the two colouring where, writing $n=2^km$ with $m$ odd, we colour $n$ red if $m\geq F(k)$ and blue if $m&lt;F(k)$ (for some sufficiently quickly growing $F$). In other words, the answer to this question is no when $k=2$. The original question is open even in the case of $\aleph_0$-many colours. This is asked by Erd\H{o}s and Galvin in \cite{ErGa91}, where they note that they do not even know whether this is possible with $k=3$. In \cite{ErGa91} they do prove that, for all $k\geq 2$, in any $k$-colouring of the integers there is a sequence such that $a_n&lt;2^{2^{O(n)}}$ for all $n$ and\[\left\{ \sum_{i\in I}a_i : \textrm{finite intervals } I\right\}\]is coloured with only two colours. This is asking about a variant of Hindman's theorem (see [532] ).




References

[ErGa91] Erd\H{o}s, Paul and Galvin, Fred, Some {R}amsey-type theorems . Discrete Math. (1991), 261--269.


\bigskip
\noindent\small{Source: \url{https://www.erdosproblems.com/948}}
\end{document}
